% TOP_PROBLEM: {"id": 3168, "title": "Hamiltonian cycles in line graphs", "category_id": 3, "category": {"id": 3, "name": "graph_theory", "display_name": "Graph Theory", "description": "Problems involving graphs, networks, and their properties.", "slug": "graph-theory", "order_index": 3, "created_at": "2026-07-31T15:26:25.671Z"}, "status": "open", "problem_number": "OPG-485", "set_id": 12}
% FIRST_POSTED: 2026-10-01
% ENTRY_KIND: statement_only
% UnsolvedMath ID 3168; source catalogue code OPG-485
% !TeX program = lualatex
% Definition and sources only; this file is not an LLM solution attempt.
\documentclass[11pt,a4paper]{article}
\usepackage[margin=25mm]{geometry}
\usepackage{fontspec}
\setmainfont{lmroman10-regular.otf}[BoldFont=lmroman10-bold.otf,
 ItalicFont=lmroman10-italic.otf,BoldItalicFont=lmroman10-bolditalic.otf]
\usepackage{amsmath,amssymb,mathtools}
\usepackage{unicode-math}
\setmathfont{latinmodern-math.otf}
\AtBeginDocument{\let\setminus\smallsetminus}
\usepackage{xurl,hyperref}
\hypersetup{colorlinks=true,urlcolor=blue}
\setcounter{secnumdepth}{0}
\setlength{\parindent}{0pt}
\setlength{\parskip}{5pt}
\setlength{\emergencystretch}{3em}
\begin{document}
\section*{Hamiltonian cycles in line graphs}\label{3168}
{\small UnsolvedMath ID 3168 · OPG-485 · Graph Theory · OpenGarden}
\subsection{Definitions and mathematical statement}
For a finite simple undirected graph $H$, the line graph
$L(H)$ has vertex set $E(H)$. Distinct edges of $H$ become
adjacent vertices of $L(H)$ exactly when they share an
endpoint. A graph $G$ is a line graph if it is isomorphic
to $L(H)$ for some finite simple graph $H$.

A graph is $4$-connected if it has at least five vertices
and remains connected after deletion of any set of at most
three vertices. A Hamilton cycle is a connected spanning
subgraph of degree two at every vertex, hence a simple
cycle visiting the entire graph.

The conjecture states that every $4$-connected line graph
has a Hamilton cycle. In terms of a root graph $H$, whenever
$L(H)$ is $4$-connected there should be a cyclic ordering
\[
 e_1,e_2,\ldots,e_m\qquad(m=|E(H)|)
\]
of all edges of $H$ such that $e_i$ and $e_{i+1}$ share
an endpoint for every index modulo $m$.

This ordering is the Hamilton cycle in $L(H)$. It need
not be an Euler tour in $H$: consecutively shared endpoints
need not give a consistent traversal of each original edge.
The hypothesis concerns vertex deletion in the line graph,
not vertex connectivity of $H$. In particular, it should
not be replaced by the different assumption that the root
graph itself is $4$-edge-connected. The root graph is useful
for expressing the target, but the conjecture is a property
of $G=L(H)$.
\subsection{Short English statement}
Does every line graph that remains connected after deletion of any three vertices contain a Hamilton cycle?
\subsection{Sources}
\begin{itemize}
\item[S1] ULAM.AI and the UnsolvedMath Contributors, supplied problems.json, record 3168 (OPG-485), OpenGarden collection. Catalogue identity and source transcription; CC BY 4.0.
\url{https://huggingface.co/datasets/ulamai/UnsolvedMath}.
\item[S2] Open Problem Garden, Hamiltonian cycles in line graphs. Original problem and discussion; the mathematical formulation below is expanded from the supplied source record.
\url{http://www.openproblemgarden.org/op/hamiltonian_cycles_in_line_graphs}.
\end{itemize}
\end{document}
